\chapter*{Résumé étendu en français}
\addstarredpart{Résumé étendu en français}
\fancyhead{}
\fancyhead[RO,LE]{\textsc{Résumé étendu en français}}
\markboth{Résumé étendu en français}{}

\section*{Introduction}

La théorie de l'Inflation constitue le socle du modèle standard de la cosmologie : une phase d'expansion accélérée de l'Univers primordial fournissant un mécanisme naturel à l'origine des fluctuations de densité observées aujourd'hui, tout en résolvant plusieurs problèmes cosmologiques (horizon, platitude et monopole). Bien que largement acceptée par la communauté, l'Inflation n'a encore reçu aucune confirmation observationnelle directe. La signature la plus prometteuse de cette phase est la composante dite des \emph{modes B} de la polarisation du fond diffus cosmologique (CMB), un motif de polarisation à divergence nulle généré par les ondes gravitationnelles primordiales prédites par l'Inflation. L'amplitude de ce signal est quantifiée par le rapport tenseur-sur-scalaire $r$, dont la mesure constituerait une preuve directe de l'existence de perturbations tensorielles primordiales de la métrique.

La détection des modes B primordiaux représente l'un des plus grands défis de la cosmologie observationnelle moderne. Le signal recherché est extrêmement faible, plusieurs ordres de grandeur sous les anisotropies de température, et il est noyé sous plusieurs sources de contamination : les avant-plans astrophysiques galactiques (poussière galactique et synchrotron), le lentillage gravitationnel qui convertit une partie des modes E en modes B, et, pour les expériences au sol, les émissions atmosphériques. Relever ce défi impose donc un contrôle rigoureux des effets systématiques instrumentaux ainsi que le développement d'algorithmes de séparation de composantes suffisamment robustes pour isoler le signal cosmologique de ses contaminants.

L'expérience QUBIC (Q\&U Bolometric Interferometer for Cosmology) propose une réponse originale à ce défi en combinant deux techniques a priori distinctes : la sensibilité des détecteurs bolométriques et le contrôle des effets systématiques propre à l'interférométrie. Cette thèse retrace le travail effectué sur cet instrument, depuis sa calibration jusqu'aux algorithmes de reconstruction de cartes qu'il rend possibles, et propose enfin une application innovante de ses capacités à l'atténuation de la contamination atmosphérique.

\section*{Contexte cosmologique}

Le modèle standard de la cosmologie, $\Lambda$CDM, décrit un Univers en expansion, homogène et isotrope à grande échelle, dominé par l'énergie noire et la matière noire froide. Le CMB, rayonnement fossile émis lors du découplage matière-rayonnement, constitue l'une des sondes les plus puissantes de ce modèle. Son spectre de température, mesuré avec une précision remarquable, ainsi que le spectre de puissance angulaire de ses anisotropies, ont permis de contraindre les paramètres cosmologiques avec une précision inégalée, notamment grâce à la mission Planck.

La polarisation du CMB se décompose en deux composantes appelées modes E et modes B. Alors que les modes E ont été mesurés avec une grande précision, les modes B primordiaux demeurent non détectés, seules des limites supérieures ayant été établies par les expériences actuelles (BICEP/Keck, POLARBEAR, SPT, ACT). La principale limitation de ces mesures provient de la contamination par les avant-plans astrophysiques, en particulier la poussière galactique, dont l'émission polarisée croît avec la fréquence, et le synchrotron, dominant aux basses fréquences. Une caractérisation précise du comportement spectral de ces avant-plans, par le biais d'une résolution en fréquence suffisante, est essentielle pour espérer isoler le signal cosmologique.

\section*{L'instrument QUBIC}

QUBIC adopte une approche instrumentale novatrice : l'interférométrie bolométrique. Plutôt que de former une image directe du ciel comme le font les imageurs classiques, l'instrument combine le signal provenant d'un réseau de 400 cornets (\emph{back-to-back horns}) pour produire un motif d'interférence, ou lobe synthétique, sur son plan focal, où sont placés des bolomètres à effet de transition supraconductrice (TES). Cette architecture combine deux atouts majeurs. Premièrement, la \emph{self-calibration}, technique héritée de la radio-interférométrie, permet de comparer les données obtenues à travers différentes paires de cornets pour caractériser et corriger les effets systématiques instrumentaux avec une précision inaccessible aux imageurs classiques. Deuxièmement, la dépendance en fréquence du lobe synthétique, dont la structure change avec la longueur d'onde, permet de reconstruire l'information spectrale directement à partir des données temporelles : c'est l'\emph{imagerie spectrale}.

Une part importante de ce travail a porté sur le système de positionnement GPS différentiel associé à la source de calibration de QUBIC, utilisée pour la self-calibration. Après une étude du bruit statistique de ce système, en particulier de sa dépendance à la distance entre antennes, l'incertitude qui en résulte a été propagée, par une méthode de Monte-Carlo, jusqu'à l'amplitude du faisceau de la source de calibration vue par l'instrument. L'incertitude obtenue s'est révélée faible, suggérant que le bruit du GPS ne devrait pas, à lui seul, limiter la précision de la self-calibration, une conclusion qu'il conviendra de confirmer une fois des données réelles disponibles sur site.

L'imagerie spectrale est ensuite présentée en détail : la position des pics secondaires du lobe synthétique variant avec la fréquence, il est possible de subdiviser les bandes physiques de l'instrument (150 et 220 GHz) en plusieurs sous-bandes reconstruites lors de l'analyse des données, sans aucune modification matérielle. Le nombre maximal de sous-bandes accessibles est fixé par la résolution spectrale de l'interféromètre, proportionnelle à l'inverse du nombre de cornets. Cette capacité, propre à la conception UWB (\emph{Ultra Wide Band}) finalement retenue pour l'instrument complet en lieu et place de la conception initiale à deux bandes séparées (\emph{Dual Bands}), est au cœur de l'ensemble des développements algorithmiques présentés dans cette thèse.

\section*{Reconstruction de cartes et séparation de composantes}

La partie centrale de cette thèse développe deux réponses complémentaires au problème de reconstruction posé par l'imagerie spectrale. La première, le \emph{Frequency Map-Making} (FMM), reconstruit plusieurs cartes de fréquence indépendantes au sein de la bande passante physique de QUBIC, à l'aide d'un solveur de type gradient conjugué préconditionné (PCG). Augmenter le nombre de sous-bandes reconstruites améliore la résolution spectrale disponible pour la séparation de composantes ultérieure, mais introduit en contrepartie un bruit anti-corrélé entre sous-bandes voisines, une limitation intrinsèque à l'imagerie spectrale. L'incorporation de données externes, notamment les cartes de l'expérience Planck aux bords du patch observé, permet par ailleurs de compenser le manque d'information caractéristique des zones de faible couverture.

La seconde méthode, le \emph{Component Map-Making} (CMM), constitue une extension plus directe du concept d'interférométrie bolométrique : plutôt que de reconstruire d'abord des cartes de fréquence, elle intègre directement une matrice de mélange dans l'opérateur de reconstruction, permettant d'effectuer simultanément la reconstruction de carte et la séparation des composantes astrophysiques. Deux variantes ont été développées : une approche \emph{Paramétrique}, efficace mais dépendante d'un modèle de spectre pour les avant-plans, et une approche \emph{Aveugle} (\emph{Blind}), qui ne fait aucune hypothèse a priori sur ce spectre. Cette dernière s'avère indispensable en présence de poussière dont le comportement spectral se décorrèle en fréquence, une situation où l'approche paramétrique introduirait un biais sur le signal cosmologique reconstruit. L'extension de CMM à un indice spectral spatialement variable et à la raie d'émission du monoxyde de carbone illustre la flexibilité du cadre développé, bien que la première de ces extensions présente encore, à ce jour, un problème de divergence en présence de bruit qui reste à comprendre.

La méthode des cross-spectres, appliquée aux cartes de fréquence issues du FMM, referme la boucle entre reconstruction et estimation des paramètres cosmologiques. Cette approche statistique, s'appuyant sur le paquet NaMaster pour l'estimation des pseudo-spectres de puissance, permet une estimation conjointe, par une analyse MCMC, du rapport tenseur-sur-scalaire et des paramètres des avant-plans. Une conséquence pratique importante en a été dégagée : le nombre de réalisations de bruit nécessaire pour obtenir une estimation non biaisée croît fortement avec le nombre de canaux de fréquence inclus dans l'ajustement, un effet relié au caractère biaisé de l'inverse d'un estimateur non biaisé de la matrice de covariance.

Une part substantielle de cette thèse a par ailleurs été consacrée au développement logiciel de l'instrument, à travers le paquet \emph{qubicsoft}, qui implémente l'ensemble du modèle de simulation et de reconstruction. Ce travail a comporté une refonte importante de son architecture, ainsi que la modernisation de sa chaîne d'installation et de ses dépendances logicielles (PyOperators, PySimulators). L'amélioration du réalisme des simulations a ensuite fait l'objet d'une étude dédiée, portant en particulier sur la convergence des données simulées en fonction du nombre de sous-bandes d'intégration et de l'ordre maximal de diffraction considéré. Un défaut de convergence inattendu, resté incompris pendant plusieurs mois, a finalement été attribué à un effet de pixellisation du ciel, les pics du lobe synthétique changeant de pixel HEALPix de manière discontinue avec la fréquence, illustration concrète des exigences de réalisme que ce type de simulation impose.

Enfin, une approche alternative au paradigme de modélisation directe (\emph{forward modelling}) commun au FMM et au CMM a été explorée à travers le développement du \emph{Neural Network Map-Making} (NNMM). Cette méthode utilise des réseaux de neurones guidés par la physique pour approximer l'inverse de l'opérateur instrumental, en factorisant celui-ci en ses différents sous-opérateurs physiques et en n'appliquant l'apprentissage automatique qu'aux composantes fondamentalement mal conditionnées, telles que la projection interférométrique. Ce cadre différentiable permet en outre d'estimer, avec leur incertitude associée, certains paramètres instrumentaux, grâce à l'inférence variationnelle mise en œuvre via le paquet Pyro. Bien qu'encore à un stade préliminaire, cette approche a été étendue avec succès du cadre du FMM à celui du CMM au cours de cette thèse.

\section*{Atténuation de la contamination atmosphérique}

La dernière partie de cette thèse s'attaque à une contamination spécifique aux expériences au sol et jusqu'ici non traitée : l'émission de l'atmosphère. L'observation de départ est que le lobe synthétique multi-pics d'un interféromètre bolométrique préserve les corrélations atmosphériques sur des échelles de temps plus longues que le lobe étroit d'un imageur classique, tout en offrant une meilleure résolution spectrale. Cette double propriété motive l'utilisation de l'imagerie spectrale de QUBIC pour séparer la contribution atmosphérique du signal cosmologique, plutôt que de la traiter comme un simple bruit à filtrer.

En étendant le formalisme du CMM Blind, l'atmosphère est modélisée comme une composante supplémentaire évoluant dans son propre référentiel local, advecté par le vent, distinct du référentiel céleste dans lequel évolue le ciel astrophysique. Cette extension a permis de démontrer, à titre de preuve de concept, la reconstruction conjointe du CMB, de l'atmosphère, de la matrice de mélange et du vecteur vent. Ce résultat constitue, à notre connaissance, la première application des capacités d'imagerie spectrale de QUBIC à l'atténuation de la contamination atmosphérique. Cette reconstruction demeure cependant fortement limitée par le coût de calcul et la lenteur de convergence du schéma d'optimisation alternée employé, ce qui a motivé le développement d'une méthode alternative, le \emph{Successive Map-Making}, qui s'affranchit de l'estimation explicite du vent au prix d'une description plus approximative.

\section*{Conclusion et perspectives}

L'ensemble des travaux présentés dans ce manuscrit contribue à faire converger, au sein d'un même cadre fondé sur l'imagerie spectrale, plusieurs approches jusque-là développées séparément : la reconstruction de cartes de fréquence, la séparation de composantes, l'inversion par réseaux de neurones, et désormais l'atténuation de la contamination atmosphérique. Cette thèse a également permis de consolider substantiellement le socle logiciel sur lequel reposent l'ensemble de ces méthodes.

Plusieurs perspectives se dégagent de ce travail. La plus récurrente est de nature algorithmique : le schéma d'optimisation alternée qui sous-tend le CMM, ainsi que son extension à l'atmosphère, souffre d'une convergence lente qui limite fortement son utilisation à une échelle plus réaliste. Le développement d'un solveur non linéaire capable d'estimer simultanément l'ensemble des paramètres, ou d'un schéma d'échantillonnage de Gibbs à la manière de Commander, apparaît comme le développement le plus susceptible de bénéficier à l'ensemble des méthodes présentées dans cette thèse. Des efforts logiciels allant dans ce sens sont d'ores et déjà en cours, notamment autour d'une réécriture des paquets PyOperators et PySimulators s'appuyant sur JAX. Sur le plan instrumental, le budget d'incertitude de la self-calibration lié au bruit du GPS devra être réévalué avec des données réelles obtenues sur site, et la divergence encore mal comprise de la reconstruction d'un indice spectral spatialement variable devra être résolue avant de pouvoir être utilisée sur des modèles de ciel réalistes. Concernant l'atmosphère, l'incorporation de mesures externes, qu'elles proviennent de stations météorologiques ou, de façon plus prometteuse, d'un système LIDAR dédié, permettrait de réduire la dépendance actuelle de la reconstruction aux seules données de l'instrument. Enfin, amener l'approche NNMM à un niveau de maturité comparable à celui du FMM et du CMM, notamment par le développement d'une étude de sensibilité dédiée, permettrait de comparer sur un pied d'égalité ces deux philosophies de reconstruction, aujourd'hui complémentaires plutôt que concurrentes.

Ces résultats confirment la prémisse centrale de cette thèse : les propriétés spécifiques de l'interférométrie bolométrique, la self-calibration et l'imagerie spectrale en particulier, constituent des atouts non seulement pour le contrôle des effets systématiques classiques, mais également pour la prise en charge de contaminants, à l'image de l'atmosphère, que les expériences au sol ont jusqu'à présent dû filtrer plutôt que modéliser.
